% Buku Ajar Pemrograman Game Menggunakan Unity
\documentclass[a4paper,12pt]{book}
\usepackage[a4paper, margin=2.5cm]{geometry}
\usepackage[utf8]{inputenc}
\usepackage[T1]{fontenc}
\usepackage[indonesian]{babel}
\usepackage{hyperref}
\usepackage{graphicx}
\usepackage{tikz}
\usetikzlibrary{positioning,arrows.meta}
\usepackage{csquotes}
\usepackage{subfiles}
\usepackage{booktabs}
\usepackage{tabularx}
\usepackage{titlesec}
\usepackage{setspace}
\usepackage{bookmark}
\usepackage{tocloft}

% Bibliografi dengan biblatex
\usepackage[backend=biber,style=authoryear,sorting=nyt,maxbibnames=99]{biblatex}
\addbibresource{references.bib}

% Spasi dan tampilan
\onehalfspacing
\setcounter{secnumdepth}{3}
\setcounter{tocdepth}{2}

% Makro utilitas untuk subfile: abaikan konten saat dibangun dari buku utama
\newcommand{\onlyinsubfile}[1]{}
% (Opsional) placeholder bila ingin memanggil bibliografi per-bab
\newcommand{\printchapterbibliography}{}

\title{Buku Ajar: Pemrograman Game Menggunakan Unity}
\author{Tim Pengajar Program Studi Teknik Informatika \\ Fakultas Teknologi Informasi \\ Universitas Bale Bandung}
\date{\today}

\begin{document}
\frontmatter
\maketitle
\tableofcontents

\mainmatter
% Bab-bab (subfiles) — setiap bab dapat dikompilasi mandiri
\subfile{chapters/00-frontmatter}
\subfile{chapters/01-pengenalan}
\subfile{chapters/02-dasar-csharp}
\subfile{chapters/03-gameobject-komponen}
\subfile{chapters/04-aset-scene}
\subfile{chapters/05-kontrol-input}
\subfile{chapters/06-fisika}
\subfile{chapters/07-animasi}
\subfile{chapters/08-ui}
\subfile{chapters/09-audio}
\subfile{chapters/10-ai-dasar}
\subfile{chapters/11-mekanik-permainan}

\backmatter
\printbibliography

\end{document}
